\section{Proofs for the sparse polynomial section}\label{app:sp}

This appendix proves the evaluation statement \cref{prop:sp:eval} and
\cref{cor:sp:qp}. The evaluation lemma \cref{lem:sp:gls} is stated for
general polytopes and approximate oracles, because \cref{sec:con} uses it for
simplices, order polytopes and implicit graphs. The margin bounds of
\cref{app:con} count nonsingular stationary points with the following form of
the refined B\'ezout inequality, which is also used in the proof of
\cref{lem:count:finite-tails}(d).

\subsection{Nonsingular common zeros}\label{app:sp:bezout}

\begin{lemma}[Nonsingular common zeros]\label{lem:sp:bezout}
Let $k\ge1$ and $D\ge1$, and let $P_1,\dots,P_k\in\mathbb C[u_1,\dots,u_k]$
have degree at most $D$. At most $D^k$ points $u\in\mathbb C^k$ satisfy
$P(u)=0$ and $\det\bigl(\partial P/\partial u\bigr)(u)\ne0$.
\end{lemma}

\begin{proof}
If some $P_l$ is a nonzero constant there are no zeros, and if $P_l=0$ the
Jacobian has a zero row everywhere; in both cases the claim holds. Otherwise
each zero set $V(P_l)$ is a hypersurface of degree at most $\deg P_l\le D$.
At a zero $u$ with invertible Jacobian, the holomorphic inverse function
theorem shows that $P$ is injective on a neighborhood of $u$, so $u$ is an
isolated point of $Z=V(P_1)\cap\dots\cap V(P_k)$, that is, an irreducible
component of $Z$ of degree one. By the refined B\'ezout inequality
\citep[Example~8.4.6]{fulton1998-intersection-theory}, the sum of the degrees
of the irreducible components of $Z$ is at most $\prod_l\deg P_l\le D^k$.
\end{proof}

The lemma counts only nonsingular zeros. Stationarity equations of a face
can have singular zeros and positive-dimensional components; these are
neither excluded nor counted. In every application, positive point growth is
what places the optimizer among the nonsingular zeros.

\subsection{Weak convex optimization on a certified patch}\label{app:sp:eval}

We use two problems of \citet[Section~2.1]{grotschel1988-geometric-algorithms-and-combinatorial-optimization}
for a convex body $\mathcal K\subseteq\R^m$, with
$S(\mathcal K,\varepsilon)$ the closed $\varepsilon$-neighborhood of
$\mathcal K$ and $S(\mathcal K,-\varepsilon)$ the set of points whose
$\varepsilon$-ball lies in $\mathcal K$. The \emph{weak separation problem}
asks, for a rational $w$ and a rational $\delta>0$, either to assert
$w\in S(\mathcal K,\delta)$ or to find a rational $\pi$ with
$\norm\pi_\infty=1$ and $\pi^{\mathsf T}x\le\pi^{\mathsf T}w+\delta$ for all
$x\in S(\mathcal K,-\delta)$. The \emph{weak optimization problem} asks, for
a rational objective $\pi$ and a rational $\varepsilon>0$, either to find a
rational $w\in S(\mathcal K,\varepsilon)$ with
$\pi^{\mathsf T}x\le\pi^{\mathsf T}w+\varepsilon$ for all
$x\in S(\mathcal K,-\varepsilon)$, or to assert that
$S(\mathcal K,-\varepsilon)$ is empty. By
\citet[Corollary~4.2.7]{grotschel1988-geometric-algorithms-and-combinatorial-optimization}, for a convex body
contained in the ball of radius $R$ about the origin and given by a weak
separation oracle, weak optimization is solved in oracle-polynomial time in
the dimension, the encoding lengths of $R$, $\pi$ and $\varepsilon$; by
their Section~4.1 conventions, a polynomial-time oracle with
polynomial-length answers makes this ordinary polynomial bit time. The
returned point need not lie in $\mathcal K$, and its objective is compared
with $S(\mathcal K,-\varepsilon)$ only; the proof below repairs both.

\begin{lemma}[Evaluation of a certified convex patch]\label{lem:sp:gls}
Let $P\subseteq\R^N$ be a nonempty polytope given by rational linear
equations and inequalities, and let rational data of total encoding length
$\Sigma$ be given as follows.
\begin{enumerate}[label=(\alph*)]
\item A point $c\in P$; a matrix $Z\in\Q^{N\times k}$, $k\ge1$, whose columns
span the direction space of the affine hull of $P$; and rationals
$r\in(0,\frac12]$, $R\ge1$ and $\zeta\ge\max\{1,\norm Z_2\}$ such that
$c+Zt\in P$ whenever $\norm t_2\le r$, and $\norm t_2\le R$ whenever
$c+Zt\in P$.
\item A function $f$, continuously differentiable on an open set containing
$P$ and convex on $P$; a norm $\nu$ on $\R^N$ equal to $\norm\cdot_2$ or
$\norm\cdot_\infty$, with dual norm $\nu_*$; and rationals
$V\ge\max\{1,\sup_P\abs f\}$ and $G\ge\max\{1,\sup_P\nu_*(\nabla f)\}$.
\item An oracle that, for a rational $y\in P$ and a rational $\beta>0$,
returns rationals $l\le f(y)\le u$ with $u-l\le\beta$ and a rational
$\hat g\in\Q^k$ with $\norm{\hat g-Z^{\mathsf T}\nabla f(y)}_\infty\le\beta$,
in time polynomial in $\Sigma$, the length of $y$, and $\log(1/\beta)$.
\item A repair map $\varrho$, computable in time polynomial in $\Sigma$ and
the length of its argument, that sends every $w\in\Q^N$ to a rational point
of $P$ with $\nu(\varrho(w)-\bar y)\le\norm{w-\bar y}_2$ for every
$\bar y\in P$.
\end{enumerate}
Then for every rational $\eta\in(0,1]$ one computes a rational $y\in P$ and
rationals $a\le U$ with
\[
 a\le\min_Pf\le f(y)\le U\le a+\eta
\]
in time polynomial in $\Sigma$ and $\log(1/\eta)$. If moreover $f$ is twice
continuously differentiable and $v^{\mathsf T}\nabla^2f(x)v\ge g\norm v_2^2$
for all $x\in P$, all $v$ in the column space of $Z$, and a rational $g>0$,
then $f$ has a unique minimizer $y^*$ on $P$, and
$\norm{y-y^*}_2\le\sqrt{2\eta/g}$.
\end{lemma}

\begin{proof}
Let $P'=\{t\in\R^k:c+Zt\in P\}$, $\tilde f(t)=f(c+Zt)$, and
\[
 \mathcal K=\{(t,\theta)\in\R^{k+1}:\ t\in P',\ \tilde f(t)\le\theta\le V+2\}.
\]
This is a compact convex set. It contains the ball of radius $r$ about
$o=(0,V+1)$: a point of this ball has $t\in P'$ and
$V+\frac12\le\theta\le V+\frac32$, while $\tilde f(t)\le V$. It lies in the
ball of radius $R_{\mathcal K}=R+2V+1$ about $o$, because
$-V\le\theta\le V+2$. Translating by $o$ gives a circumscribed body in the
sense of the cited corollary.

\emph{Weak separation.} Let $D_1=2kR$, so that
$\norm{t'-t}_1\le D_1$ for $t,t'\in P'$. Given a rational $(t,\theta)$ and
$\delta>0$: if $c+Zt$ violates a defining constraint $\alpha^{\mathsf T}x\le
\alpha_0$ of $P$ (equations being treated as two inequalities), then
$Z^{\mathsf T}\alpha\ne0$, because $\alpha^{\mathsf T}(c+Zt)=\alpha^{\mathsf
T}c\le\alpha_0$ would hold otherwise, and $(Z^{\mathsf T}\alpha,0)$ separates
strictly; if $\theta>V+2$, the vector $(0,1)$ separates strictly. Otherwise
call the oracle at $y=c+Zt$ with $\beta=\delta/(2D_1+2)$. If
$\theta<l-\beta D_1$, then for every $(t',\theta')\in\mathcal K$, convexity
gives
\[
 \theta'\ge\tilde f(t')\ge\tilde f(t)+\nabla\tilde f(t)^{\mathsf T}(t'-t)
 \ge l+\hat g^{\mathsf T}(t'-t)-\beta D_1>\theta+\hat g^{\mathsf T}(t'-t),
\]
because $\nabla\tilde f(t)=Z^{\mathsf T}\nabla f(y)$; so $(\hat g,-1)$
separates strictly. Otherwise
$\tilde f(t)-\theta\le u-l+\beta D_1\le\delta$, so either
$(t,\theta)\in\mathcal K$ or $(t,\tilde f(t))\in\mathcal K$ lies within
distance $\delta$; in both cases $(t,\theta)\in S(\mathcal K,\delta)$. A
strictly separating vector, divided by its $\infty$-norm, is a valid answer
for every $\delta$. All answers are rational of polynomial length.

\emph{Weak optimization.} Let $\Xi=G\zeta+2+(2V+1)/r$ and
$\varepsilon=\min\{r/2,\eta/(2\Xi)\}$, and apply weak optimization with
objective $(0,-1)$. Let $t^*$ minimize $\tilde f$ on $P'$, put
$z^*=(t^*,\tilde f(t^*))$ and $\lambda=\varepsilon/r\le\frac12$. The point
$z_\lambda=(1-\lambda)z^*+\lambda o$ lies in $S(\mathcal K,-\varepsilon)$:
for $\norm w\le\varepsilon$,
$z_\lambda+w=(1-\lambda)z^*+\lambda(o+w/\lambda)$ with
$\norm{w/\lambda}\le r$, a convex combination of points of $\mathcal K$.
So $S(\mathcal K,-\varepsilon)\ne\emptyset$, and the algorithm returns
$(t_o,\theta_o)\in S(\mathcal K,\varepsilon)$ with
\[
 \theta_o\le(z_\lambda)_\theta+\varepsilon
 \le\min_Pf+\lambda(V+1-\min_Pf)+\varepsilon
 \le\min_Pf+\varepsilon\Bigl(1+\frac{2V+1}r\Bigr).
\]
Put $a=\theta_o-\varepsilon(1+(2V+1)/r)$, so $a\le\min_Pf$.

\emph{Repair.} Some $(\bar t,\bar\theta)\in\mathcal K$ satisfies
$\norm{(t_o,\theta_o)-(\bar t,\bar\theta)}_2\le\varepsilon$. Let
$w=c+Zt_o$, $\bar y=c+Z\bar t\in P$ and $y=\varrho(w)\in P$. Then
$\norm{w-\bar y}_2\le\zeta\varepsilon$ and
$\nu(y-\bar y)\le\zeta\varepsilon$. Along the segment from $\bar y$ to $y$,
which lies in $P$, the mean value theorem gives
\[
 f(y)\le f(\bar y)+G\nu(y-\bar y)\le\bar\theta+G\zeta\varepsilon
 \le\theta_o+\varepsilon+G\zeta\varepsilon .
\]
Call the oracle at $y$ with $\beta=\eta/2$ and let $U=u$. Then
$f(y)\le U\le f(y)+\eta/2$ and
\[
 U-a\le\frac\eta2+\varepsilon\Bigl(G\zeta+2+\frac{2V+1}r\Bigr)\le\eta .
\]
All tolerances have polynomial encoding length, so the running time is
polynomial.

\emph{Distance.} Under the curvature hypothesis, $f$ is strictly convex on
$P$, so $y^*$ is unique. For $y\in P$ the function
$\varphi(\theta)=f(y^*+\theta(y-y^*))$ on $[0,1]$ has $\varphi'(0)\ge0$ by
first-order optimality on the convex set $P$, and
$\varphi''\ge g\norm{y-y^*}_2^2$ because $y-y^*$ lies in the column space of
$Z$. Hence $f(y)-f(y^*)\ge\frac g2\norm{y-y^*}_2^2$, and
$\norm{y-y^*}_2^2\le2(U-a)/g\le2\eta/g$.
\end{proof}

No inner-radius assumption on the location of the minimizer is used: the
known inner ball is centered at $c$, and boundary minimizers cause no loss.
The tolerance $\varepsilon$ is exponentially small when $r$ or $V^{-1}$ is,
but its encoding length is polynomial; the cost depends on $\log(1/r)$ and
$\log V$, not on their magnitudes.

\begin{proof}[Proof of \cref{prop:sp:eval}]
(a) If $\bar S=\emptyset$ the point is rational and there is nothing to do.
Otherwise apply \cref{lem:sp:gls} with $N=\abs{\bar S}$, the box $P$ of the
patch, $f=F_\gamma(\bar x_S,\cdot)$, $c$ its midpoint, $Z=I$, $\zeta=1$,
$r=\min\{\frac12,\min_{i\in\bar S}(b_i-a_i)/2\}$ and
$R=\max\{1,\sum_{i\in\bar S}(b_i-a_i)\}$. Take $\nu=\norm\cdot_2$; termwise
monomial bounds on $\widehat X$ give rationals $V\ge\max\{1,\sup\abs{F_\gamma}\}$
and $G\ge\max\{1,\sup\norm{\nabla F_\gamma}_1\}\ge\sup\norm{\nabla
F_\gamma}_2$ of polynomial length. The oracle evaluates $f$ and its
gradient exactly at rational points, and the repair map is coordinatewise
clipping to $P$, the Euclidean projection onto a box, which is nonexpansive
and fixes $P$. By \eqref{eq:sp:patch} the curvature hypothesis holds with
$g=g_0$. With $\eta=\min\{2^{-q},\frac{g_0}22^{-2q}\}$, the lemma gives
$y\in P$ and $a\le f^*\le F_\gamma(\bar x_S,y)\le a+2^{-q}$, where
$\min_Pf=f^*$ by \cref{prop:sp:closure-sound}, and $\norm{y-\hat y}_2\le
2^{-q}$. The encoding length of all data is $\poly_d(I+q)$.

(b) Let $x^\circ$ be the fallback optimizer and let
$G'\ge\max\{1,\sup_{\widehat X}\norm{\nabla F_\gamma}_1\}$ be a rational
monomial bound. Using the refinement of \cref{thm:count:fallback}, compute
for each $i\in I_C$ a rational $\tilde x_i$ with
$\abs{\tilde x_i-x_i^\circ}\le\epsilon:=2^{-q-1}/(nG')$, and let $y_i$ be
$\tilde x_i$ clipped to $[\ell_i,u_i]$; clipping does not increase the error,
since $x^\circ_i\in[\ell_i,u_i]$. For $i\in I_Z$ refine until the isolating
interval contains exactly one integer, which is $x_i^\circ$, and put
$y_i=x_i^\circ$. Then $y\in X$, $\norm{y-x^\circ}_\infty\le\epsilon$,
$\norm{y-x^\circ}_2\le\sqrt n\epsilon\le2^{-q}$, and the segment from $x^\circ$
to $y$ lies in $\widehat X$, so
$0\le F_\gamma(y)-f^*\le G'\epsilon\le2^{-q-1}$. Refine the value to an
interval of width at most $2^{-q-1}$ and let $a$ be its lower endpoint. Then
$a\le f^*\le F_\gamma(y)\le a+2^{-q}$. All refinements use $q+\poly_d(I)$
bits of precision, so the cost is $B\poly_d(I+q)$ by
\cref{thm:count:fallback}, after enlarging the base factor if necessary.

The expectation follows because the fallback output occurs with probability
at most $1/(2B)$.
\end{proof}

\subsection{The quadratic case}\label{app:sp:qp}

\begin{lemma}[Face enumeration for quadratics]\label{lem:sp:faces}
Let $F_\gamma(x)=\frac12x^{\mathsf T}Hx+(b+\gamma)^{\mathsf T}x+c_0$ be
rational on a mixed box $X$. For every integer assignment $z$ and every
partition $I_C=\mathsf L\cup\mathsf U\cup\mathsf F$, if
$H_{\mathsf F\mathsf F}$ is positive definite (or $\mathsf F=\emptyset$),
solve the linear system $\nabla_{\mathsf F}F_\gamma(\tilde x)=0$ for the point
$\tilde x$ with integer part $z$, values $\ell_{\mathsf L}$, $u_{\mathsf U}$
on the fixed sets and unknowns on $\mathsf F$, and keep $\tilde x$ if it lies
in $X$. Some kept point is a global minimizer of $F_\gamma$ on $X$. The
enumeration costs $3^{n_c}N_Z\poly(I+b)$ bit operations, where $b$ bounds the
bit length of the coefficients of $\gamma$, and every kept point is rational
of bit length $\poly(I+b)$.
\end{lemma}

\begin{proof}
Let $z^*$ be the integer part of some global minimizer. Among the global
minimizers with integer part $z^*$ choose $x^*$ whose set $\mathsf F$ of
continuous coordinates strictly inside their bounds is smallest. If
$\mathsf F=\emptyset$, $x^*$ is a kept point. Otherwise
$\nabla_{\mathsf F}F_\gamma(x^*)=0$ and $H_{\mathsf F\mathsf F}\succeq0$ by the
first- and second-order conditions in the open box. If
$H_{\mathsf F\mathsf F}v=0$ for some nonzero $v$ supported on $\mathsf F$,
then $F_\gamma(x^*+\theta v)=F_\gamma(x^*)$ for all $\theta$, because the
gradient term and the quadratic term along $v$ vanish; moving $\theta$ until
a coordinate of $\mathsf F$ reaches a bound gives a global minimizer with
integer part $z^*$ and fewer interior coordinates, a contradiction. So
$H_{\mathsf F\mathsf F}$ is positive definite and $x^*$ is the unique
solution of its system, hence kept. Every kept point is feasible, so a kept
point of least value is a global minimizer. There are $N_Z3^{n_c}$ systems;
each is a rational linear system whose solution has polynomial bit length
by Cramer's rule and Hadamard's inequality.
\end{proof}

The argument uses that a quadratic is constant along a null direction of
its Hessian at a stationary point. It does not extend to polynomials of
higher degree, for which \cref{thm:count:fallback} is used instead.

\begin{proof}[Proof of \cref{cor:sp:qp}]
Run \cref{def:sp:algorithm} with the factor $B=3^{n_c}N_Z$ of
\cref{lem:sp:faces} in the schedule and with this enumeration as the
fallback; the scalar-section count $C_{\rm tail}$ of \cref{lem:sp:tails} is
the one for $d=2$. All statements about pruning, counting, closure and the
schedule are those of the case $d=2$ of \cref{thm:sp:main}; $\kappa_3$ may be
replaced by $0$ because the Hessian is constant. When the closure test returns
a point, it is rational. When it returns $(S,\bar x_S,P,g_0)$, the problem
$\min\{F_\gamma(\bar x_S,y):y\in P\}$ is a convex quadratic program with a
positive definite Hessian, whose unique solution is the global minimizer
by \cref{prop:sp:closure-sound}; it is rational of polynomial bit length
and is computed in polynomial time \citep{kozlov1980-the-polynomial-solvability-of-convex}. On
fallback draws, \cref{lem:sp:faces} returns a rational global minimizer at
cost $B\poly(I+\log M)$, with probability at most $1/(2B)$ by
\cref{lem:sp:rare}. The work bound is that of \cref{thm:sp:main}.
\end{proof}
